\begin{titlepage}




\pagestyle{empty}


\begin{textblock}{0}(1.3,0.9)
\begin{center}
\begin{tabular}{c}
{\Huge \bf Universidade de Ciência e Tecnologia de Macau\\[10pt]
\hline\\[20pt]
{\bf \huge \sc Faculty of Innovation Engineering} \\[10pt]
{\huge Electronic and Information Engineering} \\[15pt]
\\
\\
\includegraphics[width=9cm]{Resources/must-logo.png}\\[15pt]
{\huge }

{\huge \textbf{EIE 341 Experiment Class}}\\[60pt]

{\hbox to 8truecm{\hrulefill}}\\[5pt]
{\huge  \textbf{Analog Circuit}}\\[10pt]
{\huge  \textbf{Laboratory Report}}\\
{\hbox to 8truecm{\hrulefill}}\\[50pt]

\begin{minipage}[t]{10cm}
    {\Large{\bf Estudante: \\Xiaoyang Zheng}}
\end{minipage}
\begin{minipage}[t]{5cm}\raggedleft
    {\Large{\bf Estudante N°\\ 1220034170}}\\

    {\Large{\bf Classe\\ D2}}
\end{minipage}

\\
\\
\\

{\Large Sept 2024} \\ [5pt]
\hline\\[10pt]
\bf \LARGE Coloane~~~Macao SAR
\FloatBarrier


\end{tabular}
\end{center}
\end{textblock}






%se volete metterlo all'interno del file con il resto della tesi mettete questi (altrimenti vi si sfasano le pagine:

%\thispagestyle{empty}
%\phantom{}
%\eject

% se volete aggiungere una pagina bianca (qui o da qualche parte) ripetete questi 3 comandi





\end{titlepage}